
On va
montrer que pour tout \(i\), on a :
\[
  i_*(\Fil^iR)\subset \Fil^{i+d}(K^{\bullet}(X)_{\mathbb Q}).
\]
Cela entrainera en particulier :
\[
  i_*(x)\in\Fil^{k+d}K^{\bullet}(X),
\]
et montrera donc le théorème.

La démonstration se fait par récurrence descendante sur \(i\), le cas \(i=m\) étant trivial. Soit \(y\in\Fil^iR\), et posons pour tout \(j\) :
\[
  a_j=(-1)^{j-1}(j-1)!.
\]
Alors d'après V 6.10,
\[
  \gamma^{d+i}(N,y)-a_{d+i}y\in\Fil^{i+1}R.
\]

On peut donc appliquer l'hypothèse de récurrence, supposant que :
\[
  i_*(\Fil^{i+1}R)\subset \Fil^{i+d+1}K^{\bullet}(X)_{\mathbb Q}.
\]
On a par conséquent :
\[
  i_*\bigl(\gamma^{d+i}(N,y)-a_{d+i}y\bigr)
  \in\Fil^{i+d+1}K^{\bullet}(X)_{\mathbb Q}.
\]
D'après 4.4, ceci entraine :
\[
  \gamma^{d+i}\bigl(i_*(y)\bigr)-a_{d+i}i_*(y)
  \in\Fil^{i+d+1}K^{\bullet}(X)_{\mathbb Q}.
\]
Comme \(i_*(y)\) est d'augmentation nulle, \(\gamma^{d+i}\bigl(i_*(y)\bigr)\in\Fil^{d+i}K^{\bullet}(X)\).
Comme on fait des calculs modulo torsion, on trouve donc :
\[
  i_*(y)\in\Fil^{i+d}K^{\bullet}(X)_{\mathbb Q}.
\]

\textbf{Remarque 4.6.1.} On voit que, même si 4.3 était vrai sans tensoriser par \(\mathbb Q\), la démonstration ci-dessus ne marche que modulo torsion. Pour un contre-exemple, voir XIV 4.6.

